%======================================================================================
\section{Shape: chains or compact groups}\label{sec:chains}

Inside a configuration the members are fixed, but the pair intervals among them start and end.
Between two such changes the graph has a \emph{density} (the share of the member pairs in contact)
and a \emph{diameter} (the most steps between two members); its shape is \emph{complete} (every
member in contact with every other), a \emph{tree} (the fewest edges that keep the members
connected: a chain or a star) or between (Appendix~\ref{app:graphs}). The TA-MSD pieces of
\emph{Group and solo flights} lie inside one configuration and take the shape of the graph at their
middle.

\begin{figure}[H]
\centering
\includegraphics[width=\textwidth]{chains.pdf}
\caption{Left: share of the time of each shape, by size (8+ pooled). Middle: share of the time by
diameter. Right: $H_\mathrm{loc}(\tau)$ of paragliders in triangles and chains of 3 pilots, and in
complete and tree graphs of 4 or more (flights in brackets), with 5--95\% bands.}
\label{fig:chains}
\end{figure}

\begin{reading}
\textbf{Reading.}
\begin{itemize}
  \item \textbf{Large groups are strings, not clouds.} Three pilots form a triangle \StTriangleThree\% of the
        time and a chain \StChainThree\%. Complete graphs fall from \StCompleteFour\% of the time at 4 pilots to
        \StCompleteMax\% at 8+, the mean density from \StDensityThree{} at 3 to \StDensityMax{} at 8+; at 8+ the
        diameter is 4 or more \StDiamFourMax\% of the time: pilots follow each other along the route.
  \item \textbf{The shape does not change the flight.} $H_\mathrm{loc}(1000~\mathrm{s})$ is \StHlocTriangle{} in
        triangles and \StHlocChain{} in chains; \StHlocComplete{} in complete graphs of 4+ and \StHlocTree{} in
        trees. The straighter progress of \emph{Group and solo flights} does not need every pilot in contact
        with every other.
\end{itemize}
\end{reading}
